% Part: computability
% Chapter: computability-theory
% Section: halting-problem

\documentclass[../../../include/open-logic-section]{subfiles}

\begin{document}

\olfileid[pt]{cmp}{thy}{hlt}
\olsection{O problema da parada}

Por construção, a função parcialmente computável
universal~$\fn{Un}(e,x)$ está definida se, e somente se, a computação
da função codificada por~$e$ produz um valor para a entrada~$x$. É
natural perguntar se podemos decidir quando esse é o caso. Na verdade,
não podemos. Para o modelo de computação das máquinas de Turing, isso
significa que é computacionalmente indecidível se uma dada máquina de
Turing para em uma dada entrada. Por isso, o teorema a seguir é
conhecido como a ``indecidibilidade do problema da parada''.
Apresentaremos duas demonstrações a seguir. A primeira dá continuidade
ao fio de nossa discussão anterior, enquanto a segunda é mais direta.

\begin{thm}
\ollabel{thm:halting-problem}
Seja
\[
h(e, x)  =
\begin{cases}
1 & \text{se\/ $\fn{Un}(e, x)$ está definida} \\
0 & \text{caso contrário.}
\end{cases}
\]
Então $h$ não é computável.
\end{thm}

\begin{proof}
Suponha que $h$ seja computável. Mostraremos que isso nos permitiria
definir uma função computável universal. Defina
\[
\fn{Un'}(e,x) =
\begin{cases}
\fn{Un}(e,x) & \text{se $h(e,x) = 1$} \\
0 & \text{caso contrário.}
\end{cases}
\]
Mas agora $\fn{Un'}(e, x)$ é uma função total e será computável se~$h$
o for. Por exemplo, poderíamos definir $g$ por recursão primitiva:
\begin{align*}
g(0, e, x) & \simeq 0 \\
g(y+1, e, x) & \simeq \fn{Un}(e,x);
\end{align*}
então
\[
\fn{Un'}(e,x) \simeq g(h(e,x),e,x).
\]
Como $\fn{Un'}(e,x)$ coincide com $\fn{Un}(e,x)$ sempre que esta
última está definida, $\fn{Un'}$ é universal para as funções
parcialmente computáveis que são totais. Mas isso contradiz
\olref[nou]{thm:no-univ}.
\end{proof}

\begin{proof}
Suponha que $h(e,x)$ fosse computável. Defina a função $g$ por
\[
g(x) =
\begin{cases}
  0                & \text{se $h(x,x) = 0$} \\
  \fundefined & \text{caso contrário.}
\end{cases}
\]
A função $g$ é parcialmente computável. Por exemplo, podemos defini-la
como $\umin{y}{h(x,x) = 0}$. Assim, para algum~$e$,
$g(x) \simeq \fn{Un}(e, x)$ para todo~$x$. $g$ está definida em $e$?
Se estiver, então, pela definição de~$g$, $h(e,e) = 0$ ($h$~só pode
assumir o valor~$0$ se estiver definida). Pela definição de~$h$, isso
significa que $\fn{Un}(e, e)$ é indefinida. Por nossa suposição de que
$g(x) \simeq \fn{Un}(e, x)$ para todo~$x$, temos que $g(e)$ é
indefinida, uma contradição. Por outro lado, se $g(e)$ é indefinida,
então $h(e,e) \neq 0$ e, portanto, $h(e,e) = 1$. Segue-se que
$\fn{Un}(e, e)$ está definida. Mas, como
$g(x) \simeq \fn{Un}(e, x)$, $g(e)$ também estaria definida. Novamente,
uma contradição.
\end{proof}

\begin{tagblock}{TMs}
\begin{explain}
Podemos descrever esse argumento em termos de máquinas de Turing.
Suponha que existisse uma máquina de Turing~$H$ que recebesse como
entrada uma descrição de uma máquina de Turing~$E$ e uma entrada~$x$,
e decidisse se $E$ para ou não na entrada~$x$. Poderíamos então
construir outra máquina de Turing~$G$ que recebesse uma única entrada~$x$,
executasse $H$ para decidir se a máquina $M_x$ de índice~$x$ para na
entrada~$x$, e fizesse o contrário. Em outras palavras, se $H$
informar que $M_x$ para na entrada~$x$, $G$ entra em um laço infinito;
se $H$ informar que $M_x$ não para na entrada~$x$, então $G$
simplesmente para. $G$ para quando recebe seu próprio índice como
entrada? O argumento acima mostra que ela para se, e somente se, não
para---uma contradição. Logo, nossa suposição de que existe tal máquina
de Turing~$H$ deve ser falsa.
\end{explain}
\end{tagblock}

\end{document}

